%%%%%%%% GRPO can teach agents to harm each other even at a reward cost %%%%%%%%
%
% Build with `python paper/build_paper.py`.
% Terms come from peer-reviewed work where it has one; arXiv-only coinages are not
% adopted, except Gao et al. (2026): "groupmates".
%   action-dependent / action-independent baseline   Tucker et al., ICML 2018
%   leave-one-out baseline                           Kool et al. 2019; Ahmadian et
%                                                    al., ACL 2024
%   listwise (judge)                                 Cao et al., ICML 2007; Sun et
%                                                    al., EMNLP 2023
%   difference reward                                Wolpert & Tumer 2001
%   sabotage (contests)                              Lazear 1989; Konrad 2000
%   spite                                            Hamilton 1970
%   G for group size, as in GRPO.
% "Harm to groupmates" in the introduction and less formal parts; "interference"
% ("interfering action") where one word is needed. Not "externality".
% "Group baseline": computed from the rewards of the same group (defined in Sec. 2);
% "group-mean" (GRPO) or "leave-one-out" (RLOO) where the distinction matters;
% "group-relative" only for the advantage or the GRPO-style family. American spelling.
% Evidence so far:
%   Sections 2-3, App. A-B   derived and simulated
%   Section 4.1, Table 2     tested (AFFECTED_ENVIRONMENTS.md)
%   Table 3                  measured, one seed (RESULTS.md)
%   everything else          not yet run

\documentclass{article}

\usepackage{microtype}
\usepackage{graphicx}
\usepackage{subcaption}
\usepackage{booktabs}
\usepackage{float}      % [H] keeps appendix floats in their sections
\usepackage{hyperref}

% Blind submission. For a preprint use \usepackage[preprint]{icml2027};
% for the camera-ready version use \usepackage[accepted]{icml2027}.
\usepackage{icml2027}

\usepackage{amsmath}
\usepackage{amssymb}
\usepackage{mathtools}
\usepackage{amsthm}
\usepackage[capitalize,noabbrev]{cleveref}

\theoremstyle{plain}
\newtheorem{proposition}{Proposition}
\newtheorem{corollary}[proposition]{Corollary}
\theoremstyle{definition}
\newtheorem{definition}[proposition]{Definition}

% A framed box standing in for a figure or table body that does not exist yet.
% Optional argument: height.
\newcommand{\placeholder}[1][1.3in]{%
  \fbox{\parbox[c][#1][c]{\dimexpr\linewidth-2\fboxsep-2\fboxrule\relax}{%
    \centering\small\itshape placeholder}}}

% Name of the shared-machine environment (env1/ in the repository).
\newcommand{\deadlock}{Deadlock Recovery}

\icmltitlerunning{GRPO Can Teach Agents to Harm Each Other}

\begin{document}

\twocolumn[
  \icmltitle{GRPO Can Teach Agents to Harm Each Other Even at a Reward Cost}

  % Hidden by the style file in the blind version.
  \begin{icmlauthorlist}
    \icmlauthor{Pavel Golikov}{aff}
    % other authors
  \end{icmlauthorlist}

  \icmlaffiliation{aff}{Affiliation}

  \icmlcorrespondingauthor{Pavel Golikov}{email@example.org}

  \icmlkeywords{reinforcement learning, GRPO, policy gradient, baselines,
    language model agents, AI safety}

  \vskip 0.3in
]

\printAffiliationsAndNotice{}

\begin{abstract}
\end{abstract}

%%%%%%%%%%%%%%%%%%%%%%%%%%%%%%%%%%%%%%%%%%%%%%%%%%%%%%%%%%%%%%%%%%%%%%%%%%%%%%
\section{Introduction}
\label{sec:intro}

\begin{figure}[t]
  \begin{center}
    \placeholder[1.6in]
    \caption{Schematic of the incentive created by a group baseline.
      A group of $G$ rollouts for one prompt shares a machine. One rollout
      takes an action that lowers the rewards of its groupmates and leaves its own
      reward unchanged. The group mean decreases, and the advantage of that
      rollout, $A_i = r_i - \frac{1}{G}\sum_j r_j$, increases accordingly.}
    \label{fig:schematic}
  \end{center}
\end{figure}

%%%%%%%%%%%%%%%%%%%%%%%%%%%%%%%%%%%%%%%%%%%%%%%%%%%%%%%%%%%%%%%%%%%%%%%%%%%%%%
\section{Background}
\label{sec:background}

% GRPO and its relatives, baselines and the independence they assume, notation.

%%%%%%%%%%%%%%%%%%%%%%%%%%%%%%%%%%%%%%%%%%%%%%%%%%%%%%%%%%%%%%%%%%%%%%%%%%%%%%
\section{Interference Under Group Baselines}
\label{sec:theory}

\subsection{Condition for Reinforcement}
\label{sec:condition}

\begin{table}[t]
  \caption{Expected advantage of an interfering action under different
    baselines. The action costs the acting rollout $c$ in reward and lowers
    the reward of each of $k$ groupmates by $d$, in a group of size $G$.
    Entries give the difference in expected advantage between rollouts that
    take the action and rollouts that do not; the action is reinforced when
    this difference is positive. The group-mean and leave-one-out rows are
    special cases of the identity of Gao et al.\ (2026).}
  \label{tab:gap}
  \begin{center}
    \placeholder[1.1in]
  \end{center}
  \vskip -0.1in
\end{table}

\subsection{Affected Estimators}
\label{sec:estimators}

%%%%%%%%%%%%%%%%%%%%%%%%%%%%%%%%%%%%%%%%%%%%%%%%%%%%%%%%%%%%%%%%%%%%%%%%%%%%%%
\section{Interference Channels in Practice}
\label{sec:survey}

\subsection{Shared Machines}
\label{sec:survey-machines}

\begin{table*}[t]
  \caption{Training frameworks in which rollouts of the same group share state
    under their documented default configuration. Columns: shared resource,
    tested action, reward of the acting rollout, rewards of its groupmates,
    number of trials, and information about groupmates observable during an
    episode. Frameworks that isolate rollouts by default are listed in
    \cref{tab:isolated}.}
  \label{tab:survey}
  \begin{center}
    \placeholder[1.8in]
  \end{center}
  \vskip -0.1in
\end{table*}

\subsection{Listwise Judges}
\label{sec:survey-judges}

%%%%%%%%%%%%%%%%%%%%%%%%%%%%%%%%%%%%%%%%%%%%%%%%%%%%%%%%%%%%%%%%%%%%%%%%%%%%%%
\section{Experiments}
\label{sec:experiments}

\subsection{Advantage Audit Without Training}
\label{sec:audit}

\paragraph{Method.}
\paragraph{Code generation in a shared sandbox.}
\paragraph{Listwise judging.}

\begin{figure}[t]
  \begin{center}
    \placeholder
    \caption{Advantage audit of untrained models. Mean advantage of rollouts
      that take an interfering action and of rollouts that do not, under the
      group-mean baseline and under an action-independent baseline computed
      from a second group for the same prompt on a separate machine. Error
      bars denote 95\% confidence intervals over groups.}
    \label{fig:audit}
  \end{center}
\end{figure}

\begin{figure}[t]
  \begin{center}
    \placeholder
    \caption{Effect of a single inserted sentence on a listwise judge that
      scores all members of a group in one context. Change in score for the
      rollout containing the sentence and for each of its groupmates,
      relative to the unmodified group. Error bars denote 95\% confidence
      intervals over groups.}
    \label{fig:judge-sentence}
  \end{center}
\end{figure}

\subsection{Training on a Shared Machine}
\label{sec:env1}

\paragraph{Environment.}
\paragraph{Effect of the baseline.}
\paragraph{Dependence on cost.}
\paragraph{Actions that benefit groupmates.}
\paragraph{Other estimators.}
\paragraph{Generalization beyond the group.}

% Candidate for the appendix if space is short.
\begin{table}[t]
  \caption{Untrained models in the \deadlock{} environment. Each of $G = 4$ agents runs a
    job under a common process name; terminating a process by name stops the jobs of all
    agents, whereas terminating it by process ID stops one. Columns:
    mean reward, fraction of jobs completed, fraction of rollouts that
    terminate a process by name, and fraction of rollouts that terminate a running job
    of a groupmate, on a shared machine and on isolated machines.}
  \label{tab:env1-base}
  \begin{center}
    \placeholder[1.1in]
  \end{center}
  \vskip -0.1in
\end{table}

\begin{figure*}[t]
  \begin{center}
    \placeholder[1.9in]
    \caption{Training in the \deadlock{} environment for each combination of baseline
      (group-mean or action-independent) and machine (shared or isolated),
      with all other settings held fixed. Top: fraction of rollouts that take
      the interfering action. Bottom: mean reward. Lines show means over
      seeds; shaded regions show 95\% confidence intervals.}
    \label{fig:factorial}
  \end{center}
\end{figure*}

\begin{figure}[t]
  \begin{center}
    \placeholder
    \caption{Empirical test of the condition in \cref{tab:gap}. (a)~Change in
      the frequency of the interfering action over training as a function of
      its cost $c$ to the acting rollout, under the group-mean and
      action-independent baselines. The dashed line marks the predicted
      threshold $kd = (G-1)c$. (b)~The same environment with an action that
      raises the rewards of groupmates instead of lowering them.}
    \label{fig:threshold}
  \end{center}
\end{figure}

\begin{table}[t]
  \caption{Frequency of the interfering action and mean reward in
    the \deadlock{} environment on a shared machine, before and after training with each
    estimator. Means and 95\% confidence intervals over seeds.}
  \label{tab:estimators}
  \begin{center}
    \placeholder[1.1in]
  \end{center}
  \vskip -0.1in
\end{table}

% Candidate for the appendix if space is short.
\begin{table}[t]
  \caption{Generalization beyond the training group. Frequency of the
    interfering action for the untrained model and for checkpoints trained
    with the group-mean and action-independent baselines, evaluated with
    groupmates from the same policy, with scripted agents, and with processes
    owned by a user.}
  \label{tab:persistence}
  \begin{center}
    \placeholder[0.9in]
  \end{center}
  \vskip -0.1in
\end{table}

\subsection{Training Against a Listwise Judge}
\label{sec:judge-training}

\begin{figure}[t]
  \begin{center}
    \placeholder
    \caption{Training against a listwise judge. Fraction of rollouts
      containing text addressed to the judge about other group members, and
      mean judge score, under the group-mean baseline and under an
      action-independent baseline computed from a second group scored in a
      separate context. Lines show means over seeds; shaded regions show
      95\% confidence intervals.}
    \label{fig:judge-training}
  \end{center}
\end{figure}

%%%%%%%%%%%%%%%%%%%%%%%%%%%%%%%%%%%%%%%%%%%%%%%%%%%%%%%%%%%%%%%%%%%%%%%%%%%%%%
\section{Mitigations}
\label{sec:mitigations}

\subsection{Action-Independent Baselines}
\label{sec:action-independent}

\begin{figure}[t]
  \begin{center}
    \placeholder
    \caption{Cost of the action-independent baseline when rollouts cannot
      affect one another. Reward during training under the group-mean and
      action-independent baselines with an equal number of rollouts per step.
      Lines show means over seeds; shaded regions show 95\% confidence
      intervals.}
    \label{fig:cost}
  \end{center}
\end{figure}

\subsection{Monitors}
\label{sec:monitors}

% Candidate for the appendix if space is short.
\begin{figure}[t]
  \begin{center}
    \placeholder
    \caption{Monitors that require no change to training, applied to the runs
      in \cref{fig:factorial}. (a)~Mean reward of rollouts sampled from a
      frozen copy of the initial policy and included in each group. (b)~The
      advantage audit of \cref{fig:audit}, computed at every training step.}
    \label{fig:monitors}
  \end{center}
\end{figure}

%%%%%%%%%%%%%%%%%%%%%%%%%%%%%%%%%%%%%%%%%%%%%%%%%%%%%%%%%%%%%%%%%%%%%%%%%%%%%%
\section{Related Work}
\label{sec:related}

\section{Limitations}
\label{sec:limitations}

\section{Conclusion}
\label{sec:conclusion}

% Required by the ICML 2027 template; neither counts toward the page limit.
\section*{AI Use Statement}

\section*{Impact Statement}

% Accepted version only: acknowledgements, and a link to code and data. For review,
% code goes in the supplementary material, not a URL in the paper.
% \section*{Acknowledgements}

% Enable once refs.bib exists and the text cites from it.
% \bibliography{refs}
% \bibliographystyle{icml2027}

%%%%%%%%%%%%%%%%%%%%%%%%%%%%%%%%%%%%%%%%%%%%%%%%%%%%%%%%%%%%%%%%%%%%%%%%%%%%%%
\newpage
\appendix
\onecolumn

\section{Derivations}
\label{app:derivations}

\subsection{Expected Advantage of an Action}

\subsection{Affected Estimators}

\subsection{Standard-Deviation Scaling}

\subsection{Defensive Actions}

\section{Tabular Simulations}
\label{app:sim}

\begin{figure}[H]
  \begin{center}
    \placeholder[1.6in]
    \caption{Measured versus predicted difference in expected advantage
      between rollouts that take the interfering action and rollouts that do
      not, in a tabular simulation across group sizes, numbers of affected
      groupmates, and estimators. Points on the diagonal agree with
      \cref{tab:gap}; points above zero correspond to reinforced actions.}
    \label{fig:sim-gap}
  \end{center}
\end{figure}

\begin{figure}[H]
  \begin{center}
    \placeholder[1.6in]
    \caption{Probability of the interfering action after training in a
      tabular simulation, as a function of the ratio $d/c$ of the harm per
      groupmate to the cost to the acting rollout, for three values of $k$
      ($G = 8$) and each estimator. The dashed line marks the predicted
      threshold $d/c = (G-1)/k$.}
    \label{fig:sim-threshold}
  \end{center}
\end{figure}

\begin{figure}[H]
  \begin{center}
    \placeholder[1.6in]
    \caption{Probability of the interfering action after training in a
      tabular simulation, as a function of group size, when each action
      imposes a fixed total harm $D$ divided equally among all groupmates.
      The dashed line marks the predicted boundary $G = 1 + D/c$.}
    \label{fig:sim-groupsize}
  \end{center}
\end{figure}

\begin{figure}[H]
  \begin{center}
    \placeholder[1.6in]
    \caption{Tabular simulation with a defensive action that blocks incoming
      harm at a cost. Left and center: trajectories of the probabilities of
      interference and defense under the group-mean baseline; the marker
      indicates the predicted mixed equilibrium. Right: mean reward during
      training.}
    \label{fig:sim-defense}
  \end{center}
\end{figure}

\section{Survey Details}
\label{app:survey}

\subsection{Test Procedure}

\subsection{Results by Framework}

\subsection{Frameworks With Isolated Rollouts}

\begin{table}[H]
  \caption{Training frameworks that isolate rollouts by default, with a
    separate container, virtual machine, or in-memory state per rollout, and
    the source of evidence (code or documentation).}
  \label{tab:isolated}
  \begin{center}
    \placeholder[0.9in]
  \end{center}
\end{table}

\section{The \deadlock{} Environment}
\label{app:env1}

\subsection{Job, Tools, and Reward}

\subsection{Prompts}

\subsection{Variants}

\begin{table}[H]
  \caption{Hosted models in the \deadlock{} environment without training, on shared and
    isolated machines, with and without a sentence in the tool description
    recommending termination by process ID. Columns: mean reward, fraction of
    jobs completed, fraction of termination calls made by name, and fraction
    of rollouts that terminate a running job of a groupmate.}
  \label{tab:env1-hosted}
  \begin{center}
    \placeholder[1.1in]
  \end{center}
\end{table}

\section{Training Details}
\label{app:training}

\begin{table}[H]
  \caption{Training hyperparameters and compute for each experiment: model,
    estimator, group size, groups per step, learning rate, number of steps,
    number of seeds, and GPU-hours.}
  \label{tab:recipe}
  \begin{center}
    \placeholder[0.9in]
  \end{center}
\end{table}

\section{Judge Prompts}
\label{app:judge}

\section{Evaluation of Deployed Models}
\label{app:deployed}

\begin{table}[H]
  \caption{Frequency with which each model, choosing between two actions of
    equal benefit to itself, selects the one that harms other processes on
    the machine, by ownership of those processes: rollouts of the same task,
    an unrelated user, or earlier runs of the same model. Checkpoints trained
    in the \deadlock{} environment with the group-mean and action-independent baselines
    serve as controls.}
  \label{tab:deployed}
  \begin{center}
    \placeholder[0.9in]
  \end{center}
\end{table}

\section{Additional Results}
\label{app:additional}

\end{document}
